\documentclass[12pt]{article}
\usepackage[utf8]{inputenc}
\usepackage[T1]{fontenc}
\usepackage{hyperref}
\title{Yucatec Maya Language Resources for K-12 Curriculum Development}
\author{Compiled by Claude Sonnet 5.5 (claude-sonnet-5-5) for Andr\'as Kornai}
\date{}
\begin{document}
\maketitle

\section{Introduction}
Yucatec Maya (ISO 639-3: yua; Glottocode yuca1254) is a Mayan language of the Yucatec--Lacandon group. It is spoken on the Yucatán Peninsula (Mexico) and in northern Belize. Wikipedia lists about 800,000 native speakers, including about 775,000 in Mexico in the 2020 census. Speakers mostly call it simply Maya. This is a large language with an existing, though limited, bilingual-education presence in Mexico. The resources found here are few and mostly bibliographic.

\section{Primary Educational Resources}
\begin{itemize}
\item \emph{Lengua maya: Campeche, Quintana Roo y Yucatán}, a fourth-grade textbook in Maya from the Dirección General de Educación Indígena (SEP). The first edition is from 1997 and the cited reprint is from 2013. Authors named: Víctor Manuel Canto Chan, Florinda Sosa Castilla and Carlos Humberto Contreras Ortíz. Entry in the DGEI historical collection: \url{https://dgeiib.basica.sep.gob.mx/files/fondo-editorial/acervo-historico/ah_00103.pdf}
\item Academic analysis of indigenous-language textbook coverage in Mexican primary schools. It notes that few books are published in indigenous languages for grades 3--6 and that Maya lacks clear publication criteria: \url{https://dialogossobreeducacion.cucsh.udg.mx/index.php/DSE/article/view/1572/1324}
\item Press report on SEP textbooks translated into indigenous languages (2026). Maya was not confirmed among the languages covered: \url{https://www.cronica.com.mx/nacional/2026/03/02/mario-delgado-anuncia-entrega-de-libros-de-texto-en-lenguas-originarias/}
\end{itemize}

\section{Linguistic and Reference Materials}
\begin{itemize}
\item Leiden University course page on Yucatec Maya. It names Bricker, Po'ot Yah and Dzul de Po'ot, \emph{A Dictionary of the Maya Language as Spoken in Hocab\'a, Yucat\'an} (University of Utah Press, 1998) as the main dictionary, and Blair and Vermont-Salas, \emph{Spoken Yucatec Maya} (University of Chicago, 1965--67) as an older textbook: \url{https://studiegids.universiteitleiden.nl/en/courses/41846/yucatec-maya}
\item Wikipedia article on the language: \url{https://en.wikipedia.org/wiki/Yucatec_Maya_language}
\item Glottolog entry: \url{https://glottolog.org/resource/languoid/id/yuca1254.bigmap.html}
\item Language Description Heritage entry: \url{https://ldh.clld.org/languages/yua}
\item PHOIBLE phonological inventory: \url{https://phoible.org/languages/yuca1254}
\end{itemize}

\section{Religious and Community Resources}
None found in the search results.

\section{Digital and Audio Resources}
No audio, video or software resources were found in the searches.

\section{Curriculum Development Context}
Yucatec Maya has a large speaker base. The search found state-published bilingual textbooks (DGEI/SEP) and academic reference works, but no complete K-12 curriculum. Textbook coverage in indigenous languages falls off in the upper primary grades. Secondary-level materials were not found. Further searches (CONALITEG, INALI, INDEMAYA, Belize) would be needed before curriculum work. Availability of the 1998 dictionary and the 1960s course should be checked with libraries.
\end{document}